

\chapter{Erwarteter Beitrag und Ausblick}

Die Arbeit soll eine Grundlage für die Auswahl eines \ac{RUL}-Modells schaffen, wenn im Betrieb mit unzuverlässigen Sensordaten zu rechnen ist. Dazu werden die Prognosequalität und der Aufwand der drei Modelle unter gemeinsamen Bedingungen verglichen. Vorausgesetzt wird nicht, dass MOMENT-1-small einen Vorteil hat. Sollte das Modell bei ähnlicher oder schlechterer Genauigkeit mehr Ressourcen benötigen, wäre auch dies ein wichtiges Ergebnis für die Entscheidung über seinen Einsatz.

Weitere Untersuchungen könnten prüfen, ob sich die Ergebnisse mit realen Industriedaten oder anderen \ac{TSFM} bestätigen. Auch ein gezieltes Training mit fehlerhaften Daten und eine aufwendigere Rekonstruktion fehlender Werte wären mögliche Erweiterungen. Welche davon besonders sinnvoll sind, lässt sich nach den geplanten Versuchen genauer beurteilen.
